\documentclass[10pt,a4paper]{amsart}
\usepackage[margin=19mm]{geometry}
\usepackage{amsmath,amssymb,mathtools}
\usepackage[scr=rsfs,cal=euler]{mathalfa}
\usepackage{xfrac}
\usepackage[unicode,hidelinks,bookmarks=false]{hyperref}
\newcommand{\ie}{i.e.\ }
\newcommand{\cf}{cf.\ }
\newcommand{\resp}{respectively\ }
\newcommand{\Subsection}{\subsection}
\providecommand{\nobreakdash}{\nobreak\hbox{-}}
\setlength{\emergencystretch}{2em}
\multlinegap0pt
\usepackage{bm}
\usepackage{bbm}
\usepackage{dsfont}
\usepackage{xspace}
\usepackage{cancel}
\usepackage{mathtools}
\usepackage{amsthm}
\makeatletter
\@ifundefined{Theorem}{\newtheorem{Theorem}{Theorem}}{}
\@ifundefined{Lemma}{\newtheorem{Lemma}[Theorem]{Lemma}}{}
\makeatother
\newcommand{\Q}{\mathbb{Q}}
\newcommand{\Z}{\mathbb{Z}}
\newcommand{\calN}{\mathcal{N}}
\title[Reidemeister torsion of two-bridge knots and signatures of T]{Reidemeister torsion of two-bridge knots and signatures of TQFT}
\author{Julien March\'e, Seokbeom Yoon}
\date{Source: arXiv:2511.13129v1 (CC BY 4.0)}
\begin{document}
\maketitle

\noindent\emph{Source and licence.} Reidemeister torsion of two-bridge knots and signatures of TQFT. Version arXiv:2511.13129v1 (CC BY 4.0), distributed under the Creative Commons Attribution 4.0 International licence, as recorded in the arXiv licence metadata for this exact version and shown on the abstract page https://arxiv.org/abs/2511.13129v1. Full rights notice: licenses/arxiv-2511.13129-CC-BY-4.0.txt.\par\smallskip

\noindent\textit{Editorial scope.} Counted unit: the Lemma whose source label is \texttt{Qlemma} in the article (source0.tex lines 1285--1302, source label \texttt{Qlemma}), delivered together with the article's own complete original proof of it (source0.tex lines 1303--1357, one \texttt{proof} environment of 55 lines). No mathematics is written, repaired or re-derived: statement and proof are the article's own text line for line. Required context retained uncounted: none. Every definition, display and label of those lines is retained unchanged. Omitted from this package and not redistributed: 1--1284, 1358--1523. Nothing inside a retained excerpt is omitted or altered: every retained body is delivered exactly as the article's lines are written, so no omission receipt of any kind accompanies this package. Citations: the retained excerpts carry no citation command at all, so no bibliography entry of the article is transcribed and no reference is redistributed. Reference adaptation: none was needed and none was made; every reference inside the delivered unit resolves inside it, so no unresolved reference or citation is produced. Printed slips kept literal: no printed word, symbol, hypothesis or number of the counted statement or of its proof was altered. Layout and class: the article's own class is \texttt{amsart} with packages including fontenc, graphicx, amsmath, amsfonts, amssymb, color, amsthm, hyperref, xy, tikz-cd, enumitem; the wrapper is the lane's pinned \texttt{amsart} editorial wrapper at 10pt on A4, which restates the theorem-like environments the retained text needs and adds the article's own macro definitions that the retained text needs (\texttt{\textbackslash Q}, \texttt{\textbackslash Z}, \texttt{\textbackslash calN}), with extra wrapper packages (bm, bbm, dsfont, xspace, cancel, mathtools). The delivered numbering is the unit's own. Assets: no retained span contains a figure environment, a graphics-inclusion command, a PDF-page inclusion, a table or a TikZ picture; no image file is copied into this package, the package declares no asset dependency and no image file is redistributed with it. Licence: version arXiv:2511.13129v1 of this article is distributed under the Creative Commons Attribution 4.0 International licence, as recorded in the arXiv licence metadata for this exact version and shown on the abstract page https://arxiv.org/abs/2511.13129v1. A full rights notice is delivered at licenses/arxiv-2511.13129-CC-BY-4.0.txt. Characters: every retained excerpt is pure ASCII, so the delivered engine, XeLaTeX with its default Latin Modern text font, reports no missing character.
\begin{Lemma}\label{Qlemma}
The polynomial $Q_M$ satisfies the symmetries 
$$Q_M(-U,V)=Q_M(U,V)=-Q_M(U,-V)=-Q_M(U^{-1},V^{-1})$$ 
together with the following properties: 
\begin{enumerate}
\item The terms with highest (resp. lowest) power of $V$ are $\pm V^d(U+U^{-1})^{d-1}U^{c}$ and $\pm V^{-d}(U+U^{-1})^{d-1}U^{-c}$ respectively.
%The extremal coefficients of $Q_M$ relatively to the variable $V$ are $(U+U^{-1})^{d-1}$ modulo $\Z[U^{\pm 1}]^{\times}$.
\item Its Newton polygon is the parallelogram of vertices $\pm (d-1,0)\pm(c,d)$.
\item $Q_M(\pm 1,V)=\pm 2^{d-1}(V-V^{-1})^d$.
\item $Q_M(e^u,e^v)=2^du^dH_1(v/u)$ modulo terms in $(u,v)$ of total order $<d$. 
%\item We have $\partial_VQ_M(e^u,e^v)=H_2(u,v)+\cdots$ where $H_2$ is an homogeneous polynomial of order $d-1$. 
\end{enumerate}
Similarly, the polynomials $R_M$ and $S_M$ satisfy
\begin{enumerate}
\item[(5)] $R_M(e^u,e^v)=2^du^dH_2(v/u)$  modulo terms in $(u,v)$ of total order $<d$. 
\item[(6)] $S_M(e^u,e^v)=2^bu^bH_3(v/u)$ modulo terms in $(u,v)$ of total order $<b$.  
\end{enumerate}
\end{Lemma}

\begin{proof}
%We first write $k_r-k_{r-1}=n+\alpha_r$ and check that $\alpha_r=\lfloor \frac{rc}{d}\rfloor-\lfloor\frac{(r-1)c}{d}\rfloor$ if $r<d$ and $\alpha_d=c-\lfloor\frac{(d-1)c}{d}\rfloor-1$. We observe for later use that $\alpha_r\ge 0$ and $\sum_{r=1}^d \alpha_r=c-1$. 
Before all, we provide an explicit formula for $Q_M$. Set $X=it+it^{-1},U=t,V=t^n$, and $N_k=\begin{pmatrix} i(t+t^{-1}) & 1 \\ 1 & 0 \end{pmatrix}^k(t-t^{-1})$. We compute
\begin{eqnarray*}
N_{n+\alpha}&=&\begin{pmatrix}i^{n+\alpha}t^{n+\alpha+1}-i^{n+\alpha}t^{-n-\alpha-1}& i^{n+\alpha-1}t^{n+\alpha}-i^{n+\alpha-1}t^{-n-\alpha} \\ i^{n+\alpha-1}t^{n+\alpha}-i^{n+\alpha-1}t^{-n-\alpha} & i^{n+\alpha-2}t^{n+\alpha-1} -i^{n+\alpha-2}t^{-n-\alpha+1}\end{pmatrix}\\
&=&i^{n+\alpha}\begin{pmatrix}U^{\alpha+1}V-U^{-\alpha-1}V^{-1}& i^{-1}U^{\alpha}V-i^{-1}U^{-\alpha}V^{-1}\\ i^{-1}U^{\alpha}V-i^{-1}U^{-\alpha}V^{-1}& -U^{\alpha-1}V+U^{-\alpha+1}V^{-1}\end{pmatrix}
\end{eqnarray*}
%We also compute that $-N_{n+\alpha}(-t)$ is the complex conjugate of $N_{n+\alpha}(t)$. 
If we index the elements of these matrices by signs, say $\eta,\xi\in \{\pm 1\}$, the entry $\eta,\xi$ of the matrices above are given by 
$$ \left(N_{n+\alpha}\right)_{\eta,\xi}=i^{n+\alpha-1+\frac{\eta+\xi}{2}}\Big(U^{\alpha+\frac{\eta+\xi}{2}}V-U^{-\alpha-\frac{\eta+\xi}{2}}V^{-1}\Big).$$
It follows that setting $\calN=N_{n+\alpha_1}\overline{N_{n+\alpha_{2}}}\cdots \overline{N_{n+\alpha_{d-1}}}N_{n+\alpha_d}$, its entry $\calN_{\xi_0,\xi_d}$ equals:
$$\sum_{\xi_1,\ldots,\xi_{d-1}} i^{\sum_{r=1}^d (-1)^{r+1}(n+\alpha_r-1+\frac{\xi_{r-1}+\xi_r}{2})}\prod_{r=1}^d\Big(U^{\alpha_r+\frac{\xi_{r-1}+\xi_r}{2}}V-
U^{-(\alpha_r+\frac{\xi_{r-1}+\xi_r}{2})}V^{-1}\Big)$$
We compute that the power of $i$ reduces to $n+\kappa-1+\frac{\xi_d+\xi_0}{2}$ where $\kappa=\sum (-1)^{r+1}\alpha_r$ and the product can further be expanded as follows:
$$\sum_{\xi_1,\ldots,\xi_{d-1}\in\{\pm 1\}}\sum_{\eta_1,\ldots,\eta_d\in\{\pm1\}}\Big(\prod_{r=1}^d\eta_r\Big) U^{\sum_{r=1}^d \eta_r (\alpha_r+\frac{\xi_{r-1}+\xi_r}{2})}V^{\sum_{r=1}^d \eta_r}.$$
Finally, we sum over $\xi_1,\ldots,\xi_{d-1}$ to obtain
\begin{equation}\label{PolQ}
\sum_{\eta_1,\ldots,\eta_d\in\{\pm1\}}\Big(\prod_{r=1}^d\eta_r\Big) U^{\sum_{r=1}^d \eta_r \alpha_r+\frac{\eta_1\xi_{0}+\eta_d\xi_d}{2}}\prod_{r=1}^{d-1}\Big(U^{\frac{\eta_r+\eta_{r+1}}{2}}+U^{-\frac{\eta_r+\eta_{r+1}}{2}}\Big)V^{\sum_{r=1}^d \eta_r}.
\end{equation}
To recover $Q_M(U,V)$ from this formula, we simply take $\xi_0=\xi_d=+1$ and multiply the result by the sign $i^{\kappa-1}$. It follows readily from this expression that $Q_M$ contains only even powers of $U$ and odd powers of $V$. The symmetry $Q_M(U^{-1},V^{-1})=-Q_M(U,V)$ is clear from the change $\eta_r\mapsto -\eta_r$. 

From this formula, one sees that the maximal power of $V$ in $Q_M$ is $d$ and corresponds to putting $\eta_r=1$ for all $r$: its coefficient is 
$$i^{\kappa-1} U^{\sum \alpha_r+1}(U+U^{-1})^{d-1}=i^{\kappa-1} U^{c}(U+U^{-1})^{d-1}.$$
The case of the minimal coefficient is similar and gives $i^{\kappa-1}U^{-c}(U+U^{-1})^{d-1}$.

Let us prove the second statement, we consider the linear form $L:\Z^2\to\Q$ given by $L(i,j)=i-\frac{c}{d}j$. It takes the values $\pm (d-1)$ on the sides of the parallelogram expected as the Newton polygon of $Q_M$. Hence one needs to prove that any non trivial monomial $U^iV^j$ appearing in $Q_M$ satisfies $|L(i,j)|\le d-1$. It reduces to taking a sequence $\eta_1,\ldots,\eta_d\in \{\pm 1\}$ and compute 
\begin{align*}
\begin{split}
L \Big(\sum_{r=1}^d \eta_r \alpha_r&+\frac{\eta_1+\eta_d}{2}+\sum_{r=1}^{d-1}\pm\frac{\eta_r+\eta_{r-1}}{2},\ \sum_{r=1}^d \eta_r\Big)\\
&=\sum_{r=1}^{d}\eta_r\alpha_r+\frac{\eta_1+\eta_d}{2}+\sum_{r=1}^{d-1} \pm \frac{\eta_r+\eta_{r+1}}{2}-(\sum_{r=1}^d \eta_r)\frac{c}{d}.
\end{split}
\end{align*}

Let $S$ be the set of indices $0< r<d$ for which $\eta_r\ne \eta_{r+1}$: one has 
$$\Big|\sum_{r=1}^{d-1} \pm \frac{\eta_r+\eta_{r+1}}{2}\Big|\le d-1-\#S$$
so that we are reduced to showing $|\sum \eta_r\alpha_r+\frac{\eta_1+\eta_d}{2}-(\sum\eta_r)\frac{c}{d}|\le \#S$. We do a discrete integration by parts in each sum.

On one hand: $\sum_{r=1}^d\eta_r\alpha_r=\sum_{r=1}^d\eta_r(\beta_r-\beta_{r-1})=\sum_{r=0}^{d-1}(\eta_r-\eta_{r+1})\beta_r+\eta_d(c-1)$ where $\beta_r=\lfloor \frac{rc}{d}\rfloor$ for $0\le r<d$ and $\beta_d=c-1$. 
The sequence $\eta_r-\eta_{r+1}$ vanishes if $r\notin S$ so that this sum reads $2\sum_{r\in S}\eta_r\beta_r+\eta_d(c-1)$.

On the other hand: $\sum_{r=1}^d \eta_r=\sum_{r=1}^{d-1}r(\eta_r-\eta_{r+1})+d\eta_d=2\sum_{r\in S}r\eta_r+d\eta_d$. We are then reduced to 
$$\Big|2\sum_{r\in S}\eta_r(\beta_r-r\frac{c}{d})+\frac{\eta_1-\eta_d}{2}\Big|\le \#S.$$
As $-1\le \beta_r-\frac{rc}{d}\le 0$ and the sequence $\eta_r$ is alternating, the sum of two consecutive terms does not exceed $1$ in absolute value. If $\eta_1=\eta_d$, $\#S$ is even and the result follows. When $\eta_1\ne\eta_d$ then $\#S$ is odd: either the first or the last term of the sum has a sign opposite to $\frac{\eta_1-\eta_d}{2}$: again their sum does not exceed 1 in absolute value and the argument repeats and the second point is proved. 
We also observe that in case of equality, the sequence of $\xi$s and the sequence of $\eta$s should be constant: this shows that the non trivial coefficient on the slanted sides occur only at vertices. 

Let us prove the third statement, we simply put $U=1$ in Equation \eqref{PolQ}. We also observe from the expression just above that powers of $U$ in $Q_M$ have the form $\sum_{r=1}^d \eta_r(\alpha_r+\frac{\xi_{r-1}+\xi_{r}}{2})$ which is congruent modulo 2 to $\sum_{r=1}^{d}\alpha_r+\frac{\xi_0+\xi_d}{2}=c$. Hence $Q$ has only even powers of $U$ and the statement follows. 

Finally the last three statements are clear once we have observed that writing $U=e^u, V=e^v$ one has 
\begin{align*}
i^{n+\alpha}&\begin{pmatrix}U^{\alpha+1}V-U^{-\alpha-1}V^{-1}& i^{-1}U^{\alpha}V-i^{-1}U^{-\alpha}V^{-1}\\ i^{-1}U^{\alpha}V-i^{-1}U^{-\alpha}V^{-1}& -U^{\alpha-1}V+U^{-\alpha+1}V^{-1}\end{pmatrix}\\
&=i^{n+\alpha}\begin{pmatrix} 2\sinh((\alpha+1)u+v) & \frac{2}{i}\sinh(\alpha u+v)\\ \frac{2}{i}\sinh(\alpha u+v) & -2\sinh((\alpha-1)u+v)\end{pmatrix}\\
&=2i^{n+\alpha}\begin{pmatrix} (\alpha+1)u+v & i^{-1}(\alpha u+v)\\ i^{-1}(\alpha u+v) & -(\alpha-1)u-v\end{pmatrix}+\text{higher order terms}.
\end{align*}
Taking a product of $d$ such matrices, the first order will be given by the product of the first orders. The formulas $(4,5,6)$ are direct consequences of the definition of $H_1,H_2,H_3$.  
\end{proof}

\end{document}
